\documentclass[12pt,a4paper]{article}
\usepackage[utf8]{inputenc}
\usepackage[OT4,plmath]{polski}
\usepackage{amsmath,amssymb,amsthm}
\usepackage[a4paper,margin=27mm]{geometry}
\usepackage{microtype}
\usepackage{graphicx}
\usepackage{eso-pic}
\usepackage[hidelinks,unicode]{hyperref}

\hypersetup{pdftitle={Nieliniowa Analiza Funkcjonalna}}
\setlength{\parindent}{1.5em}
\setlength{\parskip}{0pt}
\allowdisplaybreaks[1]
\theoremstyle{definition}
\newtheorem{definicja}{Definicja}
\newtheorem{przyklad}{Przykład}
\newtheorem{uwaga}{Uwaga}
\theoremstyle{plain}
\newtheorem{fakt}{Fakt}
\newtheorem{twierdzenie}{Twierdzenie}
\newtheorem{lemat}{Lemat}
\newtheorem{wniosek}{Wniosek}
\renewcommand{\proofname}{Dowód}
\newcommand{\R}{\mathbb{R}}
\newcommand{\C}{\mathbb{C}}
\newcommand{\N}{\mathbb{N}}
\newcommand{\Z}{\mathbb{Z}}

\begin{document}

% Okładka wypełnia całą pierwszą stronę A4.
\hypersetup{pageanchor=false}
\begin{titlepage}
  \thispagestyle{empty}
  \AddToShipoutPictureBG*{%
    \AtPageLowerLeft{%
      \includegraphics[width=\paperwidth,height=\paperheight]{../grafika/okladka.png}%
    }%
  }
  \null
\end{titlepage}
\hypersetup{pageanchor=true}
\pagenumbering{arabic}
\pagestyle{plain}

% =====================================================================
% TU WPISUJ CAŁĄ TREŚĆ OPRACOWANIA.
% Wszystko pozostaje w tym jednym pliku, bez rozdziałów i wykładów.
% Zapisz: Ctrl+S. Zbuduj PDF w katalogu głównym: Ctrl+Shift+B.
% =====================================================================



\end{document}
